
\newabbreviation
[type=glossaire]
{llm_1}{Grand modèle de langage (GML)} {«~Modèle de langage constitué par apprentissage profond à partir de mégadonnées.~» (tirée de \cite{Gdt_llm_0})}

\newabbreviation
[type=glossaire]
{entiteinfo}{Entité informatique}{Système, sous-système, composant logiciel ou composant matériel en lien avec l’informatique. (inspirée de~: 
\cite{iso_24765_2017, Gdt_entite_0, Gdt_informatique_0})}

\newabbreviation
[type=glossaire]
{logiciel}{Logiciel}{Ensemble de suites d’instructions interprétables par un ordinateur. (inspirée de~: \cite{iso_24765_2017, Gdt_logiciel_0, Gdt_programme_0})}

\newabbreviation
[type=glossaire]
{composantlogiciel}{Composant logiciel}{Logiciel ou partie d’un logiciel. (inspirée de~: 
\cite{iso_24765_2017, Larousse_composant_0, Gdt_composantlogiciel_0})}

\newabbreviation
[type=glossaire]
{document}{Document}{Écrit servant d’information. (inspirée de~: \cite{iso_24765_2017, Larousse_document_0})}

\newabbreviation
[type=glossaire]
{documenteninformatique}{Documentation informatique}{Ensemble de documents associé à un ensemble d’entités informatiques. (inspirée de~: 
\ccite{clements_documenting_2014}[p.~12][iso_24765_2017][][Gdt_documentation_0][])}

\newabbreviation
[type=glossaire]
{procede}{Procédé}{Ensemble structuré d'activités visant un but logiquement déterminé par des antécédents et des conséquents où des intrants donnent lieu à des extrants. (inspirée de~: \cite{Bdl_procede_0, Gdt_processus_1, Gdt_processus_0, Robert_methode_0})}

\newabbreviation
[type=glossaire]
{procedededev}{Procédé de développement}{Procédé dont le but est de créer des entités informatiques. (inspirée de~: \cite{iso_24765_2017, Larousse_developpement_0, Gdt_sdlc_0,  pmi_guidefr_2017})}

\newabbreviation
[type=glossaire]
{processus}{Processus}{Instance d'un procédé. (inspirée de~: \cite{Gdt_cycle_0, Gdt_processus_1, Gdt_processus_0, pmi_guidefr_2017})}

\newabbreviation
[type=glossaire]
{artefact}{Artefact}{Intrant tangible ou extrant tangible à un processus. (inspirée de~:  \ccite{iso_19506_2012}[p.8][Gdt_artefact_0][])}

\newabbreviation
[type=glossaire]
{phasep}{Phase d'un projet}{Regroupement d’activités logiquement interreliées découlant de la gestion d’un projet. (inspirée de~: \cite{iso_24765_2017, Larousse_phase_0})}

\newabbreviation
[type=glossaire]
{phased}{Phase d'un procédé de développement}{Regroupement d’activités logiquement interreliées découlant d’un procédé de développement. (inspirée de~: \cite{iso_24765_2017, Larousse_phase_0})}

\newabbreviation
[type=glossaire]
{realiser}{Réaliser}{Rendre tangible ce qui est abstrait. (inspirée de~: \cite{Larousse_concret_0, Larousse_tangible_0, Robert_realiser_0})}


\newabbreviation
[type=glossaire]
{dettetechnique}{Dette technique}{Passif technique résultant de choix opportunistes dont le coût de mise à niveau augmente avec le temps. (inspirée de~: \cite{avgeriou_reducing_2016, holvitie_technical_2018, Larousse_passif_0})}

\newabbreviation
[type=glossaire,sort={etablir}]
{etablir}{Établir}{Créer les conditions pour soutenir la pérennité. (inspirée de~: \cite{Robert_etablir_0, Larousse_etablir_0})}

\newabbreviation
[type=glossaire]
{instrument}{Instrument}{«~Objet fabriqué servant à exécuter [quelque chose], à faire une opération.~» (tirée de \cite{Robert_instrument_0})}

\newabbreviation
[type=glossaire]
{methode}{Méthode}{«~Ensemble de démarches raisonnées, suivies pour parvenir à un but.~» (tirée de \cite{Robert_methode_0})}

\newabbreviation
[type=glossaire]
{technique}{Technique}{«~Ensemble des applications de la science dans le domaine de la production.~» (tirée de \cite{Larousse_technique_0})}

\newabbreviation
[type=glossaire]
{organiser}{Organiser}{«~Doter d’une structure, d’une constitution déterminée, d’un mode de fonctionnement.~» (tirée de \cite{Robert_organiser_0})}

\newabbreviation
[type=glossaire]
{module}{Module}{Est un composant comprenant une interface. (inspirée de~: \ccite{clements_documenting_2014}[p.~29-31][iso_24765_2017][p.~281])}

%\newabbreviation
%[type=glossaire]
%{cohesion}{Cohésion}{Mesure les liaisons selon une logique entre les éléments d'un même composant logiciel. (inspirée de~: \ccite{iso_26514_2021}[][lano_arch_2023][p.~24,237][Larousse_cohesion_0][])}

\newabbreviation
[type=glossaire]
{couplage}{Couplage}{Mesure l’interdépendance entre des composants logiciels. (inspirée de~:  \ccite{iso_26514_2021}[][lano_arch_2023][p.~237][Larousse_coupler_0][])}

\newabbreviation
[type=glossaire]
{adaptabilite}{Adaptabilité}{La capacité d'une entité à être modifiée afin de répondre à des changements d'environnements ou de besoins. (inspirée de~:  \ccite{bass_software_2021}[p.~117][Antidote_0][][iso_9126part1_2001][p.~10][Larousse_evolution_0][])}

\newabbreviation
[type=glossaire]
{langage}{Langage}{Système pour communiquer comprenant des symboles et des règles. (inspirée de~: \cite{Gdt_langage_0, Robert_langage_0, Robert_langue_0})}

\newabbreviation
[type=glossaire]
{langageformel}{Langage formel}{Est composé d’ensembles explicites de symboles et d’un ensemble explicite de règles appelé grammaire formelle. (inspirée de~: \ccite{hunter_metalogic_1973}[p.~4][iso_24765_2017][][Gdt_lf_0][][oregan_mathematics_2020][p.~188][roggenbach_formal_2021])}%[p.~6]

\newabbreviation
[type=glossaire]
{langagenaturel}{Langage naturel}{Est un langage non formel. (inspirée de~: \cite{iso_24765_2017, Gdt_ln_0})}

\newabbreviation
[type=glossaire]
{architecturelog}{Architecture informatique}{Organisation d'entités informatiques résultant de la phase de conception. (inspirée de~: \ccite{bass_software_2021}[p.~1,3][iso_42010_2022][p.~2][meyer_agile_2014][p.~37][printz_architecture_2012])}%[p.~61]

\newabbreviation
[type=glossaire]
{architecturedescription}{Langage de description d'architecture}{Langage avec une syntaxe et une sémantique spécifiées qui décrit l’architecture d’entité informatique. (inspirée de~: \ccite{iso_42010_2022}[p.~2,20])}

\newabbreviation
[type=glossaire]
{ontologie}{Ontologie informatique}{«~Corpus structuré de concepts, qui est modélisé dans un langage permettant l’exploitation par un ordinateur des relations sémantiques ou taxonomiques établies entre ces concepts.~» (tirée de \cite{Gdt_OntInfo_1})}

\newabbreviation
[type=glossaire]
{traitementautomatiquedes}{Traitement automatique des langues (TAL)}{«~Technique d’apprentissage automatique qui permet à l’ordinateur de comprendre le langage humain.~» (tirée de \cite{Gdt_NLP_0})}

\newabbreviation
[type=glossaire]
{apprentissagepro}{Apprentissage profond (AP)}{«~Mode d’apprentissage automatique généralement effectué par un réseau de neurones artificiels composé de plusieurs couches de neurones hiérarchisées.~» (tirée de \cite{Gdt_DL_0})}

\newabbreviation
[type=glossaire]
{boitenoire}{Boîte noire}{«~Système fermé dont on ignore le fonctionnement interne, qui n’est connu que par son comportement extérieur en fonction des sollicitations qui lui sont appliquées.~» \cite{Gdt_noire_0}}

\newabbreviation
[type=glossaire]
{ellipse}{Ellipse}{«~Omission d’un ou plusieurs mots dans une phrase qui reste cependant compréhensible.~» \cite{Robert_ellipse_0}}

\newabbreviation
[type=glossaire]
{modele}{Modèle}{Représentation d’un sujet pour en permettre le traitement. (inspirée de~: \ccite{caplat_model_2008}[p.28][Robert_modele_0][][gigch_system_1991][p.135][velten_mathematical_2024][p.3])}

\newabbreviation
[type=glossaire]
{systemelogique}{Système logique}{Est composé d’un langage formel et d’un ensemble de règles (axiome, proposition, logique) permettant d’inférer. (inspirée de~: \ccite{gabbay_logical_1994}[p.217][hunter_metalogic_1973][p.7][oregan_mathematics_2020])}%[p.211]

\newabbreviation
[type=glossaire]
{machineabstra}{Machine abstraite}{Système formel comprenant les instructions nécessaires au traitement automatique de données. (inspirée de~: \cite{Wolfram_abstractmachine_0,Larousse_ordinateur_0})}

\newabbreviation
[type=glossaire]
{modeleformel}{Modèle formel}{Est composé d’un langage formel et d’un ensemble de règles et vise la représentation d’un sujet pour en permettre le traitement. (inspirée de~: \ccite{abrial_modeling_2010}[p.176][abrial_refinement_2007][][davidsson_simulation_2017][p.785])}

\newabbreviation
[type=glossaire]
{coherent}{Cohérent}{Système logique où aucune contradiction n'a encore été découverte. (inspirée de~: \ccite{Gdt_coherence_0}[][Gdt_consistance_0][][oregan_mathematics_2020][p.214])}

\newabbreviation
[type=glossaire]
{couverture}{Couverture}{Ratio de la représentation d'une chose par une autre. Cela peut être d'un modèle par un artefact, d'un sujet par un modèle ou d'une théorie par un modèle. (inspirée de~: \cite{Antidote_0, Larousse_adequat_0, Larousse_adequation_0, Larousse_complet_0, Robert_exhaustif_0})}

\newabbreviation
[type=glossaire]
{completude}{Complétude}{Un système formel est complet lorsque toutes les vérités qu'il contient peuvent être déduites des axiomes et des règles d'inférence. (inspirée de~: \ccite{Larousse_completude_0}[][oregan_mathematics_2020][p.214])}

\newabbreviation
[type=glossaire]
{comprehensible}{Compréhensible}{Qui peut être intelligible par le lectorat. (inspirée de~: \ccite{krogstie_quality_2016}[p.58-60][Larousse_comprehensible_0][][olive_conceptual_2007][p.28-30])}

\newabbreviation
[type=glossaire]
{arithmetiquedepeano}{Arithmétique de Peano}{Système formel de l'arithmétique élémentaire qui permet des raisonnements par récurrence. (inspirée de~: \ccite{delahaye_godel_2025}[p.14,66,199])}

\newabbreviation
[type=glossaire]
{outildegenie}{Outil de génie logiciel assisté par ordinateur (\textit{CASE tool}, en anglais)}{Instrument servant à réaliser ou à établir des entités. (inspirée de~: \cite{iso_24765_2017,spurr_case_1990})}

\newabbreviation
[type=glossaire]
{meta}{Méta (\textit{meta} en grec)}{«~L’élément méta‑ vient du grec meta, qui peut signifier “au milieu”, “avec” ou “après”. \elide il exprime la succession, le changement, la participation, ou un concept qui englobe un autre concept.~» (tirée de \cite{Gdtmeta_1})}

\newabbreviation
[type=glossaire]
{langagespecifique}{Langage spécifique à un domaine (\textit{Domain-specific language}, en anglais)}{Langage destiné aux utilisateurs dans un domaine spécifique dont la syntaxe ou la sémantique peut être défini formellement. (inspirée de : \ccite{Gdt_dsl_0}[][wasowski_dsl_2023][p.26-36])}

\newabbreviation
[type=glossaire]
{automate}{Automate}{Cas particulier de machine abstraite formellement défini et conçu pour en permettre l'étude. (inspirée de~: \ccite{hopcroft_introduction_2007}[p.1][kandar_introduction_2013][p.48][Gdt_automate_0][])}

\newabbreviation
[type=glossaire]
{connaissance}{Connaissance}{Une croyance justifiée qui à force de légitimités devient un fait. (inspirée de~: \cite{gingras_qu_2026})}

\newabbreviation
[type=glossaire]
{savoireninfor}{Savoir en informatique appliquée}{Ensemble des connaissances s’appliquant à l’information, aux traitements de l’information et aux entités informatiques utilisé en pratique. (inspirée de~: \cite{Gdt_1_savoir_1, Gdt_1_savoir_2})}

\newabbreviation
[type=glossaire]
{gestiondelaconnaissance}{Gestion de la connaissance}{Organiser, contrôler et rendre accessibles les représentations des connaissances au bénéfice d’une organisation. (inspirée de~: \ccite{Gdt_km_0}[][pmi_guidefr_2017][p.136])}

\newabbreviation
[type=glossaire]
{informaticien}{Informaticien}{Spécialiste du traitement automatique et rationnel de l’information. (inspirée de~: \cite{Larousse_informatique_0,Gdt_informatique_0})}

\newabbreviation
[type=glossaire]
{informaticienapp}{Informaticien appliqué}{Informaticien spécialiste dans l'information provenant de domaines autres que l'informatique. (inspiré de : \cite{Gdt_inginf_0, Gdt_swe_0, Larousse_applique_0})}

\newabbreviation
[type=glossaire]
{praticieninformatique}{Praticien informatique}{Personne qui réalise ou établit des entités informatiques. (inspirée de~: \cite{Larousse_praticien_0, Gdt_praticien_0, Gdt_informatique_0})}

\newabbreviation
[type=glossaire]
{professionnelinformatique}{Professionnel informatique}{Personne exerçant une activité rémunérée liée à l'informatique. (inspirée de~: \cite{Larousse_professionnel_0, Gdt_informatique_0})}

\newabbreviation
[type=glossaire]
{participantinformatique}{Participant informatique}{Personne participant à des processus utilisant l’informatique. (inspirée de~: \cite{Larousse_participant_0, Larousse_participer_0, Gdt_informatique_0})}

\newabbreviation
[type=glossaire]
{taxinomie}{Taxinomie}{Activité et résultat de l’étude d’un sujet pour en obtenir des règles de classification. (inspirée de~: \cite{gdt_taxinomie_0, ralph-process-2019})}

\newabbreviation
[type=glossaire]
{retroingenierie}{Rétro-ingénierie}{Consiste à étudier un groupe d’artefacts pour en déterminer le fonctionnement et la méthode de fabrication. (inspirée de~: \cite{Gdt_RetroIng_1, iso_24765_2017})}

\newabbreviation
[type=glossaire]
{norme}{Norme}{«~Document, établi par consensus et approuvé par un organisme reconnu, qui fournit, pour des usages communs et répétés, des règles \elide~» (tirée de \ccite{iso_guide2_2004}[p.12])}

\newabbreviation
[type=glossaire]
{organismedeno}{Organisme de normalisation}{«~\elide principales fonctions, en vertu de ses statuts, est la préparation, l'approbation ou l'adoption de normes \elide.~» (tirée de \ccite{iso_guide2_2004}[p.12])}

\newabbreviation
[type=glossaire]
{reglement}{Règlement}{«~Document qui contient des règles à caractère obligatoire et qui a été adopté par une autorité.~» (tirée de \ccite{iso_guide2_2004}[p.16])}

\newabbreviation
[type=glossaire]
{standard}{Standard}{Document qui contient des règles de pratiques communes ne provenant pas d’un organisme de normalisation. (inspirée de~: \cite{iso_guide2_2004, Larousse_standard_0, Gdt_standard_0})}

\newabbreviation
[type=glossaire]
{convention}{Convention}{«~Accord de volonté entre deux ou plusieurs personnes physiques ou morales, par lequel elles s'engagent à faire ou à ne pas faire quelque chose.~» (tirée de \cite{Gdt_convention_0}}

\newabbreviation
[type=glossaire]
{patrimonial}{Patrimonial}{«~Se dit de composants, logiciels ou matériels, issus d'une génération ou d'une version dépassée et qui continuent d'être utilisés, après des ajustements, en même temps que la technologie contemporaine d'une entreprise.~» (tirée de \cite{Gdt_patrimonial_0})}




